% GENERATED FILE. Edit canonical lessons, then run npm run generate:books.
% canonical-source-hash: fnv1a64:6125b387607f8489
% canonical-lessons: SA-C408-vahah, SA-C408-bharavahah, SA-C408-adhoranah, SA-C408-ankusah, SA-C408-asvasala

\chapter{Beasts of Burden}
\label{ch:sa-beasts-of-burden}
\hlchaptermodality{408}

\begin{chapteropening}
\textbf{By the end of this chapter:} \emph{I can talk about the animals that carry people and loads.}

Complete the last lesson of chapter 408: I can talk about the animals that carry people and loads.
\end{chapteropening}

\section[\texorpdfstring{\sk{वाहः} (vāhaḥ)}{वाहः (vāhaḥ)}]{\sk{वाहः} (vāhaḥ) — a draught animal}

\label{lesson:SA-C408-vahah}



\begin{quote}
Before the new one: say the Sanskrit for a goad, then the Sanskrit for harnesses.
\end{quote}

\subsection*{You'll want to know: \sk{वाहः}}
\textbf{\sk{वाहः}} — \emph{vāhaḥ} — \textquotedblleft{}a draught animal\textquotedblright{}.

\begin{quote}
\sk{ग्रामे} \sk{वृषभः} \sk{वाहः} \sk{अस्ति।}
\end{quote}

\begin{quote}
\emph{grāme vṛṣabhaḥ vāhaḥ asti.}
\end{quote}

\begin{quote}
In the village the bull is the draught animal.
\end{quote}

A masculine noun.

\begin{cousinweb}
\textbf{\sk{वाहः}} (\emph{vāhaḥ}) comes from the root √vah, \textquotedblleft{}to carry, to pull\textquotedblright{}, which also gives \emph{vahati}, \textquotedblleft{}carries\textquotedblright{}; MW gives \emph{vāha} as a draught animal: a horse, a bull or an ass (MW).
\end{cousinweb}

\begin{grammarlens}[title={What you've built}]
One more word for getting around.
\end{grammarlens}

\subsection*{Your turn}
\begin{itemize}
\raggedright
  \item \emph{Say it:} \emph{vāhaḥ}
  \item \emph{Say it:} \emph{vāhaḥ}, once more
  \item \emph{Say it:} \emph{grāme vṛṣabhaḥ vāhaḥ asti.} — In the village the bull is the draught animal.
  \item \emph{Recall:} say the Sanskrit for a goad, then the Sanskrit for harnesses, then say \emph{vāhaḥ} again
\end{itemize}

\begin{tcolorbox}[breakable,colback=teal!4,colframe=teal!35!black,title={Before you move on}]
What does \sk{वाहः} mean? (\textquotedblleft{}A draught animal\textquotedblright{}.) Say it once more.
\end{tcolorbox}
\section[\texorpdfstring{\sk{भारवाहः} (bhāravāhaḥ)}{भारवाहः (bhāravāhaḥ)}]{\sk{भारवाहः} (bhāravāhaḥ) — a porter}

\label{lesson:SA-C408-bharavahah}



\begin{quote}
Before the new one: say the Sanskrit for harnesses, then the Sanskrit for a draught animal.
\end{quote}

\subsection*{You'll want to know: \sk{भारवाहः}}
\textbf{\sk{भारवाहः}} — \emph{bhāravāhaḥ} — \textquotedblleft{}a porter\textquotedblright{}.

\begin{quote}
\sk{भारवाहः} \sk{यात्रिकस्य} \sk{भारं} \sk{वहति।}
\end{quote}

\begin{quote}
\emph{bhāravāhaḥ yātrikasya bhāraṃ vahati.}
\end{quote}

\begin{quote}
The porter carries the traveller's load.
\end{quote}

A masculine noun.

\begin{cousinweb}
\textbf{\sk{भारवाहः}} (\emph{bhāravāhaḥ}) is built from \emph{bhāra}, \textquotedblleft{}a load\textquotedblright{}, and \emph{vāha}, \textquotedblleft{}carrying\textquotedblright{}, from the root √vah; MW gives the compound as a porter, a carrier (MW).
\end{cousinweb}

\begin{grammarlens}[title={What you've built}]
One more word for getting around.
\end{grammarlens}

\subsection*{Your turn}
\begin{itemize}
\raggedright
  \item \emph{Say it:} \emph{bhāravāhaḥ}
  \item \emph{Say it:} \emph{bhāravāhaḥ}, once more
  \item \emph{Say it:} \emph{bhāravāhaḥ yātrikasya bhāraṃ vahati.} — The porter carries the traveller's load.
  \item \emph{Recall:} say the Sanskrit for harnesses, then the Sanskrit for a draught animal, then say \emph{bhāravāhaḥ} again
\end{itemize}

\begin{tcolorbox}[breakable,colback=teal!4,colframe=teal!35!black,title={Before you move on}]
What does \sk{भारवाहः} mean? (\textquotedblleft{}A porter\textquotedblright{}.) Say it once more.
\end{tcolorbox}
\section[\texorpdfstring{\sk{आधोरणः} (ādhoraṇaḥ)}{आधोरणः (ādhoraṇaḥ)}]{\sk{आधोरणः} (ādhoraṇaḥ) — a mahout; an elephant driver}

\label{lesson:SA-C408-adhoranah}



\begin{quote}
Before the new one: say the Sanskrit for a draught animal, then the Sanskrit for a porter.
\end{quote}

\subsection*{You'll want to know: \sk{आधोरणः}}
\textbf{\sk{आधोरणः}} — \emph{ādhoraṇaḥ} — \textquotedblleft{}a mahout; an elephant driver\textquotedblright{}.

\begin{quote}
\sk{आधोरणः} \sk{गजस्य} \sk{उपरि} \sk{उपविशति।}
\end{quote}

\begin{quote}
\emph{ādhoraṇaḥ gajasya upari upaviśati.}
\end{quote}

\begin{quote}
The mahout sits on top of the elephant.
\end{quote}

A masculine noun.

\begin{cousinweb}
\textbf{\sk{आधोरणः}} (\emph{ādhoraṇaḥ}) is the word MW gives for the rider or driver of an elephant; MW gives no derivation, so its origin is uncertain (MW).
\end{cousinweb}

\begin{grammarlens}[title={What you've built}]
One more word for getting around.
\end{grammarlens}

\subsection*{Your turn}
\begin{itemize}
\raggedright
  \item \emph{Say it:} \emph{ādhoraṇaḥ}
  \item \emph{Say it:} \emph{ādhoraṇaḥ}, once more
  \item \emph{Say it:} \emph{ādhoraṇaḥ gajasya upari upaviśati.} — The mahout sits on top of the elephant.
  \item \emph{Recall:} say the Sanskrit for a draught animal, then the Sanskrit for a porter, then say \emph{ādhoraṇaḥ} again
\end{itemize}

\begin{tcolorbox}[breakable,colback=teal!4,colframe=teal!35!black,title={Before you move on}]
What does \sk{आधोरणः} mean? (\textquotedblleft{}A mahout; an elephant driver\textquotedblright{}.) Say it once more.
\end{tcolorbox}
\section[\texorpdfstring{\sk{अंकुशः} (aṅkuśaḥ)}{अंकुशः (aṅkuśaḥ)}]{\sk{अंकुशः} (aṅkuśaḥ) — an elephant goad}

\label{lesson:SA-C408-ankusah}



\begin{quote}
Before the new one: say the Sanskrit for a porter, then the Sanskrit for a mahout.
\end{quote}

\subsection*{You'll want to know: \sk{अंकुशः}}
\textbf{\sk{अंकुशः}} — \emph{aṅkuśaḥ} — \textquotedblleft{}an elephant goad\textquotedblright{}.

\begin{quote}
\sk{आधोरणः} \sk{अंकुशेन} \sk{गजं} \sk{नयति।}
\end{quote}

\begin{quote}
\emph{ādhoraṇaḥ aṅkuśena gajaṃ nayati.}
\end{quote}

\begin{quote}
The mahout guides the elephant with the goad.
\end{quote}

A masculine noun.

\begin{cousinweb}
\textbf{\sk{अंकुशः}} (\emph{aṅkuśaḥ}) is the word MW gives for a hook, especially an elephant driver's hook; MW compares Greek \emph{ankistron}, \textquotedblleft{}a fish-hook\textquotedblright{}, and German \emph{Angel}, but gives no derivation within Sanskrit, so its origin is uncertain (MW).
\end{cousinweb}

\begin{grammarlens}[title={What you've built}]
One more word for getting around.
\end{grammarlens}

\subsection*{Your turn}
\begin{itemize}
\raggedright
  \item \emph{Say it:} \emph{aṅkuśaḥ}
  \item \emph{Say it:} \emph{aṅkuśaḥ}, once more
  \item \emph{Say it:} \emph{ādhoraṇaḥ aṅkuśena gajaṃ nayati.} — The mahout guides the elephant with the goad.
  \item \emph{Recall:} say the Sanskrit for a porter, then the Sanskrit for a mahout, then say \emph{aṅkuśaḥ} again
\end{itemize}

\begin{tcolorbox}[breakable,colback=teal!4,colframe=teal!35!black,title={Before you move on}]
What does \sk{अंकुशः} mean? (\textquotedblleft{}An elephant goad\textquotedblright{}.) Say it once more.
\end{tcolorbox}
\section[\texorpdfstring{\sk{अश्वशाला} (aśvaśālā)}{अश्वशाला (aśvaśālā)}]{\sk{अश्वशाला} (aśvaśālā) — a stable}

\label{lesson:SA-C408-asvasala}



\begin{quote}
Before the new one: say the Sanskrit for a mahout, then the Sanskrit for an elephant goad.
\end{quote}

\subsection*{You'll want to know: \sk{अश्वशाला}}
\textbf{\sk{अश्वशाला}} — \emph{aśvaśālā} — \textquotedblleft{}a stable\textquotedblright{}.

\begin{quote}
\sk{रात्रौ} \sk{अश्वः} \sk{अश्वशालायां} \sk{तिष्ठति।}
\end{quote}

\begin{quote}
\emph{rātrau aśvaḥ aśvaśālāyāṃ tiṣṭhati.}
\end{quote}

\begin{quote}
At night the horse stands in the stable.
\end{quote}

A feminine noun.

\begin{cousinweb}
\textbf{\sk{अश्वशाला}} (\emph{aśvaśālā}) is built from \emph{aśva}, \textquotedblleft{}a horse\textquotedblright{}, and \emph{śālā}, \textquotedblleft{}a hall, a house\textquotedblright{}; MW gives the compound as a stable for horses (MW).
\end{cousinweb}

\begin{grammarlens}[title={What you've built}]
One more word for getting around.
\end{grammarlens}

\subsection*{Your turn}
\begin{itemize}
\raggedright
  \item \emph{Say it:} \emph{aśvaśālā}
  \item \emph{Say it:} \emph{aśvaśālā}, once more
  \item \emph{Say it:} \emph{rātrau aśvaḥ aśvaśālāyāṃ tiṣṭhati.} — At night the horse stands in the stable.
  \item \emph{Recall:} say the Sanskrit for a mahout, then the Sanskrit for an elephant goad, then say \emph{aśvaśālā} again
\end{itemize}

\begin{tcolorbox}[breakable,colback=teal!4,colframe=teal!35!black,title={Before you move on}]
What does \sk{अश्वशाला} mean? (\textquotedblleft{}A stable\textquotedblright{}.) Say it once more.
\end{tcolorbox}
